% !TeX root = ../数值分析与科学计算笔记.tex
% 由主文档通过 \input 载入；请编译主文档。
\setcounter{chapter}{5} % 保留原书章号。
\chapter{线性方程组的迭代方法}
\label{chap:na-iterative}

\section{矩阵分裂与收敛的谱条件}
对于大型稀疏系统，直接分解的填充和存储代价可能很大。迭代法利用矩阵乘向量和简单子问题逐步逼近解。设 $A=M-N$，$M$ 可逆，令
\begin{equation}
 x_{k+1}=M^{-1}Nx_k+M^{-1}b=:Bx_k+c.
 \label{eq:na-stationary}
\end{equation}
实现中求解关于 $M$ 的方程，不必显式形成 $M^{-1}$。误差满足 $e_{k+1}=Be_k$。

\begin{theorem}[定常迭代的充要条件]
对任意初值，式~\eqref{eq:na-stationary} 收敛到唯一不动点，当且仅当 $\rho(B)<1$，其中谱半径 $\rho(B)=\max_{\lambda\in\sigma(B)}|\lambda|$，$\sigma(B)$ 为复特征值集合。
\end{theorem}
\begin{proof}
若对任意初始误差有 $B^ke_0\to0$，则对每个复特征向量亦有 $\lambda^kv\to0$，故 $|\lambda|<1$。反之，对复 Jordan 标准形，每个块的 $k$ 次幂由有限个 $k$ 的多项式乘 $\lambda^{k-j}$ 构成；当 $|\lambda|<1$ 时均趋于零，零特征值块最终为零。因此 $B^k\to0$，且 $I-B$ 可逆，不动点唯一。
\end{proof}

容易检查的充分条件为某诱导范数下 $\|B\|<1$。它不是必要条件；$\rho(B)<1$ 时某些范数下仍可能 $\|B\|>1$，初期误差也不必单调下降。

\section{Jacobi、Gauss--Seidel 与松弛}
写 $A=D+L+U$，其中 $D$ 为对角部分，$L,U$ 分别是严格下、上三角部分。本章采用此加号约定。若所有 $a_{ii}\ne0$，Jacobi 迭代为
\[
 x_i^{(k+1)}=\frac{b_i-\sum_{j\ne i}a_{ij}x_j^{(k)}}{a_{ii}},
 \qquad B_J=-D^{-1}(L+U).
\]
Gauss--Seidel 迭代在一轮内立即使用已更新分量：
\[
 x_i^{(k+1)}=\frac{b_i-\sum_{j<i}a_{ij}x_j^{(k+1)}-\sum_{j>i}a_{ij}x_j^{(k)}}{a_{ii}},
 \qquad B_G=-(D+L)^{-1}U.
\]
若 $A$ 按行严格对角占优，即 $|a_{ii}|>\sum_{j\ne i}|a_{ij}|$，则 Jacobi 迭代在无穷范数下为压缩映射。Gauss--Seidel 也收敛：令旧误差最大值为 $E$，$q=\max_i\sum_{j\ne i}|a_{ij}|/|a_{ii}|<1$；按行归纳，新误差每个分量都不超过 $qE$，因为本轮已更新的误差至多为 $E$，剩余旧误差也至多为 $E$。

逐分量把新 Gauss--Seidel 值与旧值以权 $\omega$ 混合，得到逐次超松弛法（successive over-relaxation，SOR）：
\[
 (D+\omega L)x_{k+1}=\omega b-\bigl(\omega U+(\omega-1)D\bigr)x_k.
\]
对称正定矩阵在 $0<\omega<2$ 时 SOR 收敛；这一标准结果可由逐坐标二次能量下降证明，本书不展开其完整谱分析。对一般矩阵，$0<\omega<2$ 本身不是充分条件，也不存在与矩阵无关的最佳参数。

\begin{example}[一个可直接比较的系统]
对 $A=\begin{pmatrix}4&1\\1&3\end{pmatrix}$，Jacobi 矩阵为 $\begin{pmatrix}0&-1/4\\-1/3&0\end{pmatrix}$，谱半径为 $1/\sqrt{12}$；Gauss--Seidel 矩阵为 $\begin{pmatrix}0&-1/4\\0&1/12\end{pmatrix}$，谱半径为 $1/12$。二者都收敛，且后者在此例具有更小的渐近因子；该比较不能无条件推广到所有矩阵。
\end{example}

\section{最速下降与共轭梯度}
对 SPD 矩阵 $A$，求解 $Ax=b$ 等价于最小化
\[
 \phi(x)=\tfrac12x^{\mathsf T}Ax-b^{\mathsf T}x,
 \qquad \phi(x)-\phi(x_*)=\tfrac12\|x-x_*\|_A^2,
\]
其中 $\|z\|_A=(z^{\mathsf T}Az)^{1/2}$ 为能量范数。残差 $r_k=b-Ax_k$ 是负梯度。沿 $r_k$ 精确线搜索得最速下降步
\[
 \alpha_k=\frac{r_k^{\mathsf T}r_k}{r_k^{\mathsf T}Ar_k},\qquad
 x_{k+1}=x_k+\alpha_kr_k.
\]
当残差非零时分母为正。条件数很大时不同方向曲率差异显著，最速下降可能反复折返，收敛缓慢。

\emph{共轭梯度法（conjugate gradient method，CG）}使用互相 $A$ 正交的搜索方向。给定 $x_0$，初始化 $r_0=b-Ax_0$、$p_0=r_0$；若 $r_0=0$ 直接返回。其精确算术递推为
\begin{align}
 \alpha_k&=\frac{r_k^{\mathsf T}r_k}{p_k^{\mathsf T}Ap_k},&
 x_{k+1}&=x_k+\alpha_kp_k,\notag\\
 r_{k+1}&=r_k-\alpha_kAp_k,&
 \beta_k&=\frac{r_{k+1}^{\mathsf T}r_{k+1}}{r_k^{\mathsf T}r_k},\notag\\
 p_{k+1}&=r_{k+1}+\beta_kp_k.
 \label{eq:na-cg}
\end{align}
若新残差为零，应在计算下一步前停止。每轮一次矩阵向量乘法及 $O(n)$ 次向量操作；对稀疏矩阵，乘法代价与非零元数同阶。

\begin{theorem}[CG 的投影性质与误差界]
\label{thm:na-cg}
在精确算术下，SPD 系统的 CG 在尚未到达解时产生 $A$ 共轭的非零搜索方向。第 $k$ 步在 $x_0+\mathcal K_k$ 中最小化能量误差，其中
\[
 \mathcal K_k=\operatorname{span}\{r_0,Ar_0,\ldots,A^{k-1}r_0\}.
\]
故它至多 $n$ 步终止。若 $\kappa=\kappa_2(A)>1$，则
\[
 \|x_k-x_*\|_A\le
 2\left(\frac{\sqrt\kappa-1}{\sqrt\kappa+1}\right)^k\|x_0-x_*\|_A.
\]
当 $\kappa=1$ 时非精确初值一步求解。
\end{theorem}
\begin{proof}[证明思路与误差界的依据]
将式~\eqref{eq:na-cg} 代入内积关系，可归纳得到不同残差正交、不同方向 $A$ 正交，以及两组向量生成相应 Krylov 空间。沿共轭方向的精确极小化不会破坏先前方向上的最优性，因此 $r_k$ 与 $\mathcal K_k$ 正交，这正是仿射空间中的极小化条件。非零共轭方向线性无关，故至多有 $n$ 个。由极小化性质，任取次数不超过 $k$ 且 $p(0)=1$ 的多项式，均有
\[
 \|x_k-x_*\|_A\le\max_{\lambda\in\sigma(A)}|p(\lambda)|\,\|x_0-x_*\|_A.
\]
把第~\ref{sec:na-chebyshev} 节的 Chebyshev 多项式从 $[-1,1]$ 映到 $[\lambda_{\min},\lambda_{\max}]$，并在零点归一化，即得所述上界。这里交代投影与多项式界的关键逻辑，递推正交关系的逐项归纳略去。
\end{proof}

有限精度中，共轭性与正交性可能逐渐丢失，不能把“至多 $n$ 步”当作实际软件的保证。递推残差也可能偏离真实残差，应在必要时重新计算 $b-Ax_k$。CG 的正定假设不可随意取消；非对称系统通常需要其他 Krylov 方法。

\section{预条件与可解释的停止}
\emph{预条件（preconditioning）}通过容易求解的近似矩阵改善迭代问题。对 SPD 预条件矩阵 $M=CC^{\mathsf T}$，等价变换为
\[
 C^{-1}AC^{-\mathsf T}y=C^{-1}b,\qquad x=C^{-\mathsf T}y.
\]
变换后的矩阵仍 SPD，适用于 CG。实现中用 $Mz=r$ 的求解代替显式形成 $C^{-1}$。Jacobi 预条件取 $M=D$，对 SPD 系统对角元为正。有效性取决于谱分布及每步成本，不仅是某个条件数数字。

迭代停止常用 $\|r_k\|\le\varepsilon_{\mathrm{abs}}+\varepsilon_{\mathrm{rel}}\|b\|$，配合最大迭代次数和停滞检测。该条件给出残差保证，前向误差需结合第~\ref{chap:na-direct} 章的条件数界。零右端项不能只用相对阈值；预条件残差和原问题残差也应明确区分。

\section{习题}
\begin{enumerate}[label=\arabic*.]
\item\label{exr:na-iter-jacobi} 对 $A=\begin{pmatrix}2&1\\1&2\end{pmatrix}$ 写出 Jacobi 和 Gauss--Seidel 迭代矩阵，求谱半径。
\item\label{exr:na-iter-nonnormal} 令 $B=\begin{pmatrix}1/2&10\\0&1/2\end{pmatrix}$。计算 $B^k$，说明 $\rho(B)<1$ 不保证每步误差范数下降。
\item\label{exr:na-iter-cg} 对 $A=\operatorname{diag}(1,2)$、$b=(1,1)^{\mathsf T}$、$x_0=0$，手算 CG 的前两步。
\item\label{exr:na-iter-energy} 证明 SPD 系统中 $\|x-x_*\|_A^2=r^{\mathsf T}A^{-1}r$，其中 $r=b-Ax$。
\end{enumerate}

